\documentclass[11pt]{article}
\usepackage[margin=1in]{geometry}
\usepackage{amsmath,amssymb,amsthm}
\usepackage{mathtools}
\usepackage{hyperref}
\usepackage{booktabs}

\newtheorem{theorem}{Theorem}[section]
\newtheorem{proposition}[theorem]{Proposition}
\newtheorem{lemma}[theorem]{Lemma}
\newtheorem{remark}[theorem]{Remark}

\newcommand{\Q}{\mathbf Q}
\newcommand{\F}{\mathbf F}
\newcommand{\BSD}{\operatorname{BSD}}

\title{An Explicit Non-Semistable Quadratic-Twist Family\\
Satisfying the Strong Birch--Swinnerton--Dyer Conjecture}
\author{Neo.K}
\date{13 August 2026}

\begin{document}
\maketitle

\begin{abstract}
Let
\[
E/\Q:\qquad y^2=x^3+x^2+8x-16,
\]
the elliptic curve with LMFDB label \(696.e1\) (Cremona label \(696b1\)).
Its conductor is \(696=2^3\cdot3\cdot29\), so \(E\) is not semistable.
Put \(f_2(x)=x^3+x^2+8x-16\), and let \(\mathcal P\) be the set of primes
\(q\) satisfying
\[
q\equiv1\pmod{24},\qquad
\left(\frac q{29}\right)=1,\qquad
f_2(x)\text{ irreducible in }\F_q[x].
\]
By combining the \(2\)-primary quadratic-twist theorem collected by
Banwait--Huang, Skinner's ordinary/multiplicative \(p\)-part theorem, the
additive-twist argument of Burungale--Skinner--Tian--Wan, and the
arbitrary-reduction theorem of Fouquet--Wan, we derive that every
\(E^{(q)}\), \(q\in\mathcal P\), satisfies the strong BSD conjecture.
Chebotarev gives \(\delta(\mathcal P)=1/24\).
\end{abstract}

\paragraph{Status.}
This note records a derived consequence of the cited theorem statements and
does not assert priority over all published or unpublished literature.

\section{Main theorem}

\begin{theorem}\label{thm:main}
Let
\[
E:\quad y^2=x^3+x^2+8x-16
\]
and
\[
\mathcal P=\left\{
q\text{ prime}:q\equiv1\pmod{24},\
\left(\frac q{29}\right)=1,\
x^3+x^2+8x-16\text{ irreducible mod }q
\right\}.
\]
Then \(\mathcal P\) has natural prime density \(1/24\), and for every
\(q\in\mathcal P\),
\[
L(E^{(q)},1)\ne0
\qquad\text{and}\qquad
\BSD(E^{(q)})
\]
holds.
\end{theorem}

\section{Arithmetic of the base curve}

The base curve has
\[
N_E=696=2^3\cdot3\cdot29,\qquad
\Delta_E=-2^{11}\cdot3\cdot29.
\]
It has rank \(0\), trivial rational torsion, singleton rational isogeny
class, Manin constant \(1\), and local reduction
\[
2:\ II^*\text{ additive},\qquad
3:\ I_1\text{ split multiplicative},\qquad
29:\ I_1\text{ nonsplit multiplicative},
\]
with \(v_3(\Delta_E)=v_{29}(\Delta_E)=1\).
LMFDB records maximal residual image at every odd prime.  The full BSD
formula for the base curve is covered by Creutz--Miller since
\(N_E<5000\) and the analytic rank is \(0\).
A modular-symbol computation gives
\[
L(E,1)/\Omega_E=1.
\]

The \(2\)-division cubic is
\[
f_2(x)=x^3+x^2+8x-16,
\]
which is irreducible over \(\Q\), with
\[
\operatorname{disc}(f_2)=-11136=-2^7\cdot3\cdot29.
\]
Hence its Galois closure \(L/\Q\) has group \(S_3\) and quadratic resolvent
\[
F_0=\Q(\sqrt{-174}).
\]

\section{Support primes}

\begin{lemma}\label{lem:split}
If \(q\in\mathcal P\), then \(2,3,29\) split in \(\Q(\sqrt q)\).
\end{lemma}

\begin{proof}
The congruence \(q\equiv1\pmod{24}\) gives splitting at \(2\) and \(3\),
while \((q/29)=1\) gives splitting at \(29\).
\end{proof}

\begin{lemma}\label{lem:ordinary}
Every \(q\in\mathcal P\) is a good ordinary prime for \(E\).
\end{lemma}

\begin{proof}
Irreducibility of \(f_2\bmod q\) gives a \(3\)-cycle on \(E[2]\).
Its trace in \(\mathrm{GL}_2(\F_2)\) is \(1\), so \(a_q(E)\) is odd.
For \(q\ge5\), supersingularity would imply \(q\mid a_q(E)\); the Hasse
bound then forces \(a_q(E)=0\), contradiction.
\end{proof}

\section{Chebotarev density}

\begin{proposition}\label{prop:density}
The set \(\mathcal P\) has natural prime density \(1/24\).
\end{proposition}

\begin{proof}
Put \(K=\Q(\zeta_{24},\sqrt{29})\). Then \([K:\Q]=16\), and
\[
F_0=\Q(\sqrt{-174})\subset K
\]
because \(\sqrt{-174}=\sqrt{-6}\sqrt{29}\) and
\(\Q(\sqrt{-6})\subset\Q(\zeta_{24})\).
Since \(K/\Q\) is abelian and \(F_0\) is the unique nontrivial proper
normal subfield of the \(S_3\)-extension \(L/\Q\),
\[
L\cap K=F_0,\qquad [LK:\Q]=48.
\]
The congruence and Legendre conditions specify the identity on \(K\), while
irreducibility of \(f_2\bmod q\) specifies the \(3\)-cycle class on \(L\).
These restrictions agree on \(F_0\).  The class has two elements in a group
of order \(48\), so Chebotarev gives \(2/48=1/24\).
\end{proof}

\section{The \(2\)-part}

\begin{proposition}\label{prop:2part}
For every \(q\in\mathcal P\),
\[
L(E^{(q)},1)\ne0,\qquad \BSD(E^{(q)},2).
\]
\end{proposition}

\begin{proof}
Apply Banwait--Huang, Theorem 2.14, in the branch
\(E(\Q)[2]=0\), \(\Delta_E<0\).  The base curve is optimal of rank \(0\),
has odd Manin constant and known \(\BSD(E,2)\), and
\[
v_2(L(E,1)/\Omega_E)=0.
\]
For \(d=q\), Lemma~\ref{lem:split} gives splitting of every prime dividing
\(N_E\), and irreducibility of \(f_2\bmod q\) is the required inertness in
the cubic \(2\)-division field.
\end{proof}

\section{Odd primes}

Fix \(q\in\mathcal P\) and put \(A=E^{(q)}\).
Quadratic twisting preserves absolute irreducibility of residual
representations, so \(A[p]\) is absolutely irreducible for every odd \(p\).

\begin{lemma}\label{lem:qpart}
\(\BSD(A,q)\) holds.
\end{lemma}

\begin{proof}
The prime \(q\) is good ordinary for \(E\).  Use the additive-twist clause
of Burungale--Skinner--Tian--Wan exactly as in the proof of
Banwait--Huang Proposition 2.9(1).  Banwait--Huang Remark 2.10 states that
semistability is used there only to produce the ramified-prime condition.
Here \(\ell=29\) supplies it directly:
\(v_{29}(\Delta_E)=1\), so \(E[q]\) is ramified at \(29\), and
\(29\nmid\operatorname{disc}\Q(\sqrt q)\).
\end{proof}

\begin{lemma}\label{lem:ordinaryp}
If \(p\) is an odd good ordinary prime for \(A\), then \(\BSD(A,p)\).
\end{lemma}

\begin{proof}
Apply Skinner's Theorem C.  The representation \(A[p]\) is irreducible, and
the multiplicative prime \(29\ne p\) is a ramified auxiliary prime since
\(v_{29}(\Delta_A)=1\).
\end{proof}

\begin{lemma}\label{lem:mult}
\(\BSD(A,3)\) and \(\BSD(A,29)\) hold.
\end{lemma}

\begin{proof}
Skinner's Theorem C allows \(p\ge3\).  At \(p=3\), use \(29\) as the
distinct ramified multiplicative prime; at \(p=29\), use \(3\).
Both discriminant valuations equal \(1\).
\end{proof}

\begin{lemma}\label{lem:ss}
If \(p\) is an odd good supersingular prime for \(A\), then \(\BSD(A,p)\).
\end{lemma}

\begin{proof}
Apply Fouquet--Wan Theorem 1.7.  Absolute irreducibility is already known.
At a good supersingular prime \(p\ge5\), the local residual representation
is irreducible, so the forbidden semisimplified character sum cannot occur.
At \(\ell=29\), nonsplit multiplicative reduction supplies the nontrivial
unramified quadratic Steinberg twist, while
\(p\nmid v_{29}(\Delta_A)=1\) gives residual ramification.  Thus
Fouquet--Wan Corollary 1.10 gives the \(p\)-part.

The period in that corollary is the modular-form period.  The isogeny class
of \(A\) is singleton because the same is true of \(E\) and quadratic
twisting preserves rational isogenies.  Hence \(A\) is optimal.  By
Česnavičius--Neururer--Saha, prime divisors of its Manin constant are
supported at additive reduction primes.  Those are \(2\) and \(q\), whereas
the present \(p\) is good; hence \(p\) does not affect the period comparison.
\end{proof}

\begin{proof}[Proof of Theorem~\ref{thm:main}]
Proposition~\ref{prop:density} gives the density.  Proposition~\ref{prop:2part}
gives nonvanishing and the \(2\)-part of BSD.  Every odd prime is either
\(q\), one of \(3,29\), good ordinary, or good supersingular; the four lemmas
above give the corresponding \(p\)-parts.  Gross--Zagier--Kolyvagin and the
local-to-global form of the strong BSD formula in analytic rank \(0\) then
give \(\BSD(A)\).
\end{proof}

\section*{Remark on status}
This note is intended as a source-audited derived corollary.  No claim of
priority is made without a separate literature and expert-referee audit.

\begin{thebibliography}{99}

\bibitem{BH}
B.~S. Banwait and X.~Huang,
\emph{On the identification of elliptic curves that admit infinitely many
twists satisfying the Birch--Swinnerton--Dyer conjecture},
arXiv:2601.16044v3 (2026).

\bibitem{BSTW}
A.~A. Burungale, C.~Skinner, Y.~Tian and X.~Wan,
\emph{Zeta elements for elliptic curves and applications},
arXiv:2409.01350v2 (2024).

\bibitem{CM}
B.~Creutz and R.~L. Miller,
\emph{Second isogeny descents and the Birch and Swinnerton--Dyer
conjectural formula},
J. Algebra 372 (2012), 673--701.

\bibitem{CNS}
K.~Česnavičius, M.~Neururer and A.~Saha,
\emph{The Manin constant and the modular degree},
J. Eur. Math. Soc. 26 (2024), 573--637.

\bibitem{FW}
O.~Fouquet and X.~Wan,
\emph{The Iwasawa Main Conjecture for universal families of modular motives},
arXiv:2107.13726v3.

\bibitem{Skinner}
C.~Skinner,
\emph{Multiplicative reduction and the cyclotomic main conjecture for
\(\mathrm{GL}_2\)},
arXiv:1407.1093.

\bibitem{LMFDB}
The L-functions and Modular Forms Database,
elliptic curve \(696.e1\) (Cremona \(696b1\)).

\end{thebibliography}

\end{document}
